\documentclass[10pt]{article}
\usepackage[letterpaper,margin=0.78in]{geometry}
\usepackage{amsmath,amssymb,amsthm,mathtools}
\usepackage{graphicx}
\usepackage{microtype}
\usepackage{xcolor}
\usepackage[colorlinks=true,linkcolor=blue!55!black,citecolor=blue!55!black,urlcolor=blue!55!black]{hyperref}

\newtheorem{theorem}{Theorem}
\newtheorem{lemma}[theorem]{Lemma}
\newtheorem{corollary}[theorem]{Corollary}
\newcommand{\Sph}{\mathbb S}
\newcommand{\R}{\mathbb R}
\newcommand{\Ip}{\mathcal I}

\title{The Generalized Busemann Intersection Question in Dimension Two}
\author{Candidate complete fixed-dimension resolution for arXiv:2208.03882}
\date{17 August 2026}

\begin{document}
\maketitle

\begin{center}
\fbox{\parbox{0.92\linewidth}{\textbf{Status: candidate substantial partial result, likely valid.}
Question~1 of the source has an affirmative answer in dimension $n=2$ for
every real $0<p<2$.  For $p\ne1$, equality holds exactly for centered
ellipses; at the Parseval value $p=1$, every origin-symmetric star body is an
equality case, as the source notes.  The higher-dimensional question remains
open.  Novelty is plausible but not certified.}}
\end{center}

\section{Source question and normalized form}

The source asks whether two sharp affine Fourier inequalities hold for all
origin-symmetric star bodies, with the direction reversing at $p=n/2$
\cite{Yaskin}:

\begin{center}
\includegraphics[width=0.96\linewidth]{figures/open_question_crop.png}
\end{center}

Let $d\sigma$ be normalized measure on $\Sph^{n-1}$.  The source introduces
the normalized spherical Fourier operator $I_p$, with $I_p1=1$, and shows
that for $0<p<n/2$ the upper inequality is equivalent to
\begin{equation}
 \|I_pf\|_{n/p}\leq \|f\|_{n/(n-p)},
 \qquad f=\rho_K^{\,n-p}.
 \label{eq:normform}
\end{equation}
The lower inequality for $n/2<p<n$ is the reverse norm inequality with the
same exponents.  Indeed,
\[
 \int_{\Sph^{n-1}}f^{n/(n-p)}\,d\sigma
 =\frac{\operatorname{vol}_n(K)}{\kappa_n},
\]
and the source constant is chosen so that balls give equality.  We prove
both statements when $n=2$.

\section{Idea of the proof}

\subsection*{Proof Intuition}

An even function on $\Sph^1$ is a function on the projective circle, hence
has period $\pi$.  For $0<p<1$, the integral kernel of $I_p$ is
$|\cos(\theta-\phi)|^{-p}$.  A quarter-turn of the output changes this to
$|\sin(\theta-\phi)|^{-p}$; doubling angles then identifies the projective
circle with the ordinary unit circle and this kernel with its chordal
Hardy--Littlewood--Sobolev kernel.  Lieb's sharp one-dimensional HLS theorem
is exactly the norm estimate \eqref{eq:normform}.  Its conformal extremizers
are negative powers of positive quadratic forms on the projective circle,
which are precisely the radial powers of centered ellipses.

For $1<p<2$, put $q=2-p<1$.  The spherical Fourier multipliers give
$I_qI_p=\mathrm{Id}$ on even functions.  Applying the already proved HLS
bound for $I_q$ to $I_pf$ reverses the inequality and settles the second
half.

\section{The projective-circle HLS identity}

Write $u_\theta=(\cos\theta,\sin\theta)$ and regard an even function as a
$\pi$-periodic function of $\theta$.  Let
\[
 d\mu(\theta)=\frac{d\theta}{\pi}\quad (0\leq\theta<\pi),
 \qquad
 a_p=\frac{\sqrt\pi\,\Gamma(1-p/2)}{\Gamma((1-p)/2)}.
\]
The source's integral formula for $I_p$ specializes, for $0<p<1$, to
\begin{equation}
 (I_pf)(\theta+\pi/2)
 =a_p\int_0^\pi |\sin(\theta-\phi)|^{-p}f(\phi)\,d\mu(\phi).
 \label{eq:kernel}
\end{equation}
The normalization is exact because the beta integral gives
\begin{equation}
 \int_0^\pi |\sin t|^{-p}\,d\mu(t)
 =\frac{\Gamma((1-p)/2)}{\sqrt\pi\,\Gamma(1-p/2)}=a_p^{-1}.
 \label{eq:mass}
\end{equation}

We use the following normalized form of the sharp HLS theorem.

\begin{lemma}[Sharp HLS on the projective circle]\label{lem:hls}
Let $0<p<1$ and $s=2/(2-p)$.  For $F,G\in L^s([0,\pi),\mu)$,
\begin{equation}
 a_p\left|\int_0^\pi\!\int_0^\pi
 \frac{F(\phi)G(\theta)}{|\sin(\theta-\phi)|^p}
 \,d\mu(\phi)d\mu(\theta)\right|
 \leq \|F\|_s\|G\|_s.
 \label{eq:hls}
\end{equation}
For nonzero nonnegative $F,G$, equality holds exactly when they are
proportional and, for some $0\leq r<1$ and $\beta\in\R$, have the form
\begin{equation}
 F(\theta)=A
 \bigl(1-2r\cos(2\theta-\beta)+r^2\bigr)^{-(2-p)/2}.
 \label{eq:extremizer}
\end{equation}
\end{lemma}

\begin{proof}
Set $x=2\theta$ and $y=2\phi$.  Then normalized measure becomes normalized
measure on the usual circle and
\[
 2|\sin(\theta-\phi)|=|e^{ix}-e^{iy}|.
\]
Thus \eqref{eq:hls} is the diagonal sharp HLS inequality in dimension one,
transported from $\R$ to the circle by stereographic projection.  Lieb's
sharp theorem gives the constant and all extremizers \cite{Lieb}; evaluating
at $F=G=1$ and using \eqref{eq:mass} yields the displayed normalization.
The conformal extremizers on the circle become exactly
\eqref{eq:extremizer} after undoing angle doubling.
\end{proof}

\begin{corollary}\label{cor:operator}
For $0<p<1$ and every even $f\in L^{2/(2-p)}(\Sph^1)$,
\begin{equation}
 \|I_pf\|_{2/p}\leq\|f\|_{2/(2-p)}.
 \label{eq:operator}
\end{equation}
For nonzero nonnegative $f$, equality holds exactly for the functions in
\eqref{eq:extremizer}.
\end{corollary}

\begin{proof}
The dual exponent of $2/(2-p)$ is $2/p$.  Apply Lemma~\ref{lem:hls}, use
duality, and then use \eqref{eq:kernel}; the quarter-turn preserves every
$L^r$ norm.  The equality statement is inherited from the bilinear theorem.
\end{proof}

\section{Resolution in dimension two}

\begin{theorem}\label{thm:main}
Question~1 of the source is true when $n=2$.

\begin{enumerate}
\item If $0<p<1$, inequality (6) holds for every origin-symmetric star body
$K\subset\R^2$, with equality exactly for centered ellipses.
\item At $p=1$, the Parseval identity gives equality for every
origin-symmetric star body.
\item If $1<p<2$, inequality (7) holds on its natural Fourier domain: if
$I_p(\rho_K^{2-p})\in L^{2/p}$, the displayed integral is finite and the
sharp lower bound holds; otherwise its extended value is $+\infty$.  Among
finite cases, equality holds exactly for centered ellipses.
\end{enumerate}
\end{theorem}

\begin{proof}
First suppose $0<p<1$ and set $f=\rho_K^{2-p}$.  Then $f$ is positive,
continuous, and even, while
\[
 \int_{\Sph^1} f^{2/(2-p)}\,d\sigma
 =\int_{\Sph^1}\rho_K^2\,d\sigma
 =\frac{\operatorname{area}(K)}{\pi}.
\]
Corollary~\ref{cor:operator} is therefore precisely the normalized form
\eqref{eq:normform} of source inequality (6).

It remains to translate equality.  The denominator in
\eqref{eq:extremizer} is a positive quadratic form in $u_\theta$:
\[
 1-2r\cos(2\theta-\beta)+r^2
 =u_\theta^TQ_{r,\beta}u_\theta,
\]
where
\[
 Q_{r,\beta}=(1+r^2)I-2r
 \begin{pmatrix}\cos\beta&\sin\beta\\
                 \sin\beta&-\cos\beta\end{pmatrix}.
\]
Its eigenvalues are $(1-r)^2$ and $(1+r)^2$, so it is positive definite.
Consequently equality means
\[
 \rho_K(\theta)=C\bigl(u_\theta^TQ_{r,\beta}u_\theta\bigr)^{-1/2},
\]
which is exactly the radial function of a centered ellipse.  Conversely,
every centered ellipse has this form and gives equality.

At $p=1$, the source's spherical Parseval identity applies and all bodies
are equality cases.

Finally let $1<p<2$, set $q=2-p$, and initially take smooth even $f$.  On
spherical harmonics of degree $2j$, the normalized Fourier multiplier of
$I_p$ is
\[
 (-1)^j\frac{\Gamma(1-p/2)\Gamma(j+p/2)}
 {\Gamma(p/2)\Gamma(j+1-p/2)}.
\]
The multiplier for $I_q$ is its reciprocal, hence $I_qI_p f=f$.  Applying
\eqref{eq:operator} with parameter $q$ to $h=I_pf$ gives
\[
 \|f\|_{2/(2-p)}=\|I_qh\|_{2/(2-p)}
 \leq\|h\|_{2/p}=\|I_pf\|_{2/p},
\]
which is the normalized form of source inequality (7).  The same argument
holds distributionally whenever $I_pf\in L^{2/p}$; standard approximation
gives the stated natural-domain formulation.

If equality holds, the HLS classification for $q$ says
$h=C(u_\theta^TQu_\theta)^{-p/2}$.  The HLS Euler--Lagrange equality, or
direct conformal covariance, then gives
$f=I_qh=C'(u_\theta^TQu_\theta)^{-(2-p)/2}$, so $K$ is a centered ellipse.
The converse also follows from the source's linear invariance and its
ellipsoid equality lemma.
\end{proof}

\section{Scope and higher-dimensional obstruction}

The proof is special to dimension two.  On the projective circle, the
orthogonality set of a point is another point, so a quarter-turn converts
the generalized cosine kernel into the diagonal HLS singularity.  In
$\mathbb{RP}^{n-1}$ for $n>2$, the orthogonality set is a hypersurface, not
the diagonal, and this reduction disappears.

The later stochastic-geometry theorem of Adamczak, Paouris, Pivovarov, and
Simanjuntak proves the $0<p<1$ portion corresponding to
$n/p\in\mathbb N$ after translating their negative $L_p$-intersection-body
parameter \cite{APPS}.  It does not cover continuous $p$, while the circle
argument does.  Attempts to interpolate those discrete parameters, to use
the ordinary higher-dimensional HLS theorem, and to force an axisymmetric
counterexample did not resolve $n\geq3$; these attempts are recorded with
the packet.

\begin{thebibliography}{9}
\bibitem{Yaskin}
V. Yaskin,
\emph{On a generalization of Busemann's intersection inequality},
J. Funct. Anal. \textbf{287} (2024), 110561;
\href{https://arxiv.org/abs/2208.03882}{arXiv:2208.03882}.

\bibitem{Lieb}
E. H. Lieb,
\emph{Sharp constants in the Hardy--Littlewood--Sobolev and related
inequalities}, Ann. of Math. (2) \textbf{118} (1983), 349--374.

\bibitem{APPS}
R. Adamczak, G. Paouris, P. Pivovarov, and P. Simanjuntak,
\emph{From intersection bodies to dual centroid bodies: a stochastic
approach to isoperimetry},
J. Eur. Math. Soc. (2024);
\href{https://arxiv.org/abs/2211.16263}{arXiv:2211.16263}.
\end{thebibliography}

\end{document}
